% Generated draft: requires mathematical and source review.
\documentclass[11pt]{article}
\usepackage[margin=25mm]{geometry}
\usepackage{fontspec}
\setmainfont{Latin Modern Roman}
\usepackage{amsmath,amssymb,mathtools,unicode-math}
\setmathfont{Latin Modern Math}
\usepackage{enumitem,needspace,xurl,hyperref}
\setlength{\parindent}{0pt}
\setlength{\parskip}{5pt}
\setlength{\emergencystretch}{3em}
\newcommand{\N}{\mathbb{N}}
\newcommand{\Z}{\mathbb{Z}}
\newcommand{\Q}{\mathbb{Q}}
\newcommand{\R}{\mathbb{R}}
\newcommand{\C}{\mathbb{C}}
\newcommand{\F}{\mathbb{F}}
\newcommand{\A}{\mathbb{A}}
\newcommand{\PP}{\mathbb{P}}
\newcommand{\E}{\mathbb{E}}
\newcommand{\eps}{\varepsilon}
\newcommand{\dd}{\,\mathrm{d}}
\begin{document}
\section*{UnsolvedMath ID: 10400021}
\textbf{Source problem code:} catalogue entry
\subsection{Definitions and mathematical statement}
Let \(L\subset S^3\) be a hyperbolic link, meaning that its complement \(X_L=S^3\setminus L\) carries the complete finite-volume hyperbolic structure intended by the source. Write \(\operatorname{vol}(X_L)\) for its hyperbolic volume and \(\operatorname{CS}(X_L)\) for its Chern--Simons invariant, with normalization and lift still to be specified. Let \(J_N(L)\), \(N\ge1\), denote exactly the colored-Jones-derived complex sequence used in the source; its normalization and evaluation parameter are not given in the record.

Conjecture 1.21, attributed to H. Murakami, J. Murakami, M. Okamoto, T. Takata and Y. Yokota [297], states that for every such hyperbolic link, with an appropriate choice of branch of the logarithm,
\[
 2\pi i\,\lim_{N\to\infty}\frac{\log J_N(L)}{N}
 =\operatorname{CS}(X_L)+i\operatorname{vol}(X_L).
\]
It states moreover that
\[
 \lim_{N\to\infty}\frac{J_{N+1}(L)}{J_N(L)}
 =\exp\!\left(\frac{\operatorname{CS}(X_L)+i\operatorname{vol}(X_L)}{2\pi i}\right).
 \tag{7}
\]
These are two assertions, not an inference from the first limit to the ratio limit. Their meaning requires a specification of \(J_N\), eventual nonvanishing where logs and ratios are used, a compatible logarithm-branch prescription, and the normalization and representative of the possibly period-valued Chern--Simons invariant. The displayed order \(\operatorname{CS}+i\operatorname{vol}\), factor \(2\pi i\), and exponential denominator are preserved exactly from the supplied statement; no alternative complex-volume convention has been substituted.
\subsection{Short English statement}
For every hyperbolic link, the specified colored-Jones sequence is conjectured to have both a logarithmic asymptotic and a consecutive-term ratio determined by Chern--Simons invariant plus \(i\) times hyperbolic volume. The supplied partial-progress assessment says this is proved for some knots and open in general, without listing the knots or supplying those proofs.
\subsection{Sources}
\begin{itemize}
\item Local UnsolvedMath registry, permanent ID \(10400021\), exact source code \texttt{catalogue entry}; title truncated to ``Conjecture 1.21 --- (H.'' The statement credits H. Murakami, J. Murakami, M. Okamoto, T. Takata and Y. Yokota, source reference [297]. The record does not include [297]'s bibliography entry.
\item Ohtsuki, \emph{Problems on invariants of knots and 3-manifolds} (2002), Conjecture 1.21, PDF page 24: \url{https://msp.org/gtm/2002/04/gtm-2002-04-024s.pdf}. Printed bold headings distinguish 231 numbered items from surrounding references to them.
\item Supplied extraction: pdf-font-structure; accessed 2026-07-29. Source status NEEDS\_REVIEW presents an open problem and calls for a current primary-literature audit. Registry status partially\_solved and classification PARTIAL-PROGRESS call it the Complex Volume Conjecture, report results for some knots, and say it is open in general; no further primary references are supplied.
\item Dataset attribution: Ulam AI / UnsolvedMath Contributors, CC BY 4.0. Source rights NEEDS\_REVIEW: publicly accessible author/publisher copy, redistribution terms requiring review.
\end{itemize}
\end{document}
